\section{Actual characteristic-\(\ell\) curves and generic ordinary Hecke support}

The source's generic ordinary determinant problem permits an arbitrary
tuple of distinct residual characters.  This section identifies the
actual curve map and its fixed-weight specialization without using
the single-character \texttt{Fr.Frame}.  The arithmetic containment
of the Hecke ideal remains an explicit hypothesis.

Let \(\mathcal O\) be the coefficient DVR, with finite residue field
\(k\) of characteristic \(\ell\).  Keep the actual cubic ring
\(S\), normalization and local conditions of AC-G49--AC-G58, in the
proved good-prime range.  In particular \(\ell\ne2,3,269\).
Write \(H=G_K\), \(K=\mathbb Q(\omega)\),
\(\gamma=\psi\kappa_\ell^e\), and
\[
 \overline r_\chi|_H
   =\operatorname{diag}(\overline\chi^{\,j}\overline\gamma,
                       1,
                       \overline\chi^{\,j}\overline\gamma^{-1}),
 \qquad j\in\{1,2\}.
 \tag{G61.1}
\]
An actual curve is a local continuous map
\(S\longrightarrow V\), where \(V\) is a complete DVR of
characteristic \(\ell\) with finite residue field \(k'\supset k\).

\begin{lemma}[AC-G61-pair: the actual curve belongs to the generic determinant functor]
\label{AC-G61-pair}
Put \(A=k+\mathfrak m_V\subset V\), using the canonical finite
coefficient fields.  The actual curve factors continuously through
\(A\), and \(A\) is a complete Noetherian local
\(\mathcal O\)-algebra with residue field \(k\).

Choose a source \texttt{HidaSetting} and \texttt{Lev} over a CM
field \(L\supset K\), with its actual splitting and unramifiedness
hypotheses retained.  Suppose the restricted representation descends
to its frame group and that its three normalized residual characters
are distinct.  Let \(\mathfrak r\) be the source
\texttt{W.Resid} with that residual tuple, the corresponding local
residual characters, and \(\texttt{th}=1\) on the local
\(\ell\)-primary torsion subgroup.  Use the actual away-inertia
killing subgroups and the ordered flags of AC-G57-integral.
Then the normalized representation
\[
             \rho=r_\chi|_{G_L}\,\overline\kappa_\ell^{-1}
 \tag{G61.2}
\]
and its ordered local characters define an actual
\texttt{Pair} \((\mathfrak r,A)\).
Consequently any source
\texttt{UniversalOrdinaryDeterminantRing} \(B\) for
\(\mathfrak r\) has the uniquely specified map
\[
       f:B\longrightarrow A\hookrightarrow V
 \tag{G61.3}
\]
which carries its universal determinant and local characters to
those of this curve.  Existence of \(B\), when needed, uses the
source's explicitly retained \texttt{ProModularRingExists} input.
No assertion about a Hecke ideal is used to construct \(f\).
\end{lemma}
\begin{proof}
\proofstep{The coefficient ring.}
Hensel lifting the simple roots of \(X^{|k'|}-X\) gives the unique
copy of \(k'\) in \(V\).  A uniformizer \(t\) then gives
\(V=k'[[t]]\): successive subtraction of residue coefficients
constructs the expansion, and completeness and separatedness prove
existence and uniqueness.  Thus
\(A=k+t k'[[t]]\).  If \(1,b_2,\ldots,b_f\) is a
\(k\)-basis of \(k'\), then
\[
 A=k[[t]]\cdot1+\sum_{i=2}^f k[[t]]\cdot(tb_i).
 \tag{G61.4}
\]
It is a finite module over the Noetherian ring \(k[[t]]\), hence
Noetherian.  An element is a unit precisely when its constant term
is nonzero, so its maximal ideal is \(tV\), with residue field
\(k\).  Its powers are \(t^nV\), and the expansion proves
completeness for this topology.  Every image of an element of \(S\)
has residue in \(k\), since the curve is local and respects the
coefficient map.  Hence its image lies in \(A\).  The topologies
on \(A\) and on its image in \(V\) agree, giving continuity.

\proofstep{Polarization and residual data.}
The actual semidirect representation has the checked symmetric
conjugate-dual form of multiplier one on \(H\).  After the scalar
normalization in (G61.2), that form has multiplier
\(\overline\kappa_\ell^{-2}\), exactly
\(\epsilon^{1-n}\) for the source's \(n=3\).
If \(P\) is its symmetric invertible form matrix, then
\[
 \rho(g)^tP\rho(g^c)=\overline\kappa_\ell(g)^{-2}P.
 \tag{G61.5}
\]
For the source involution on the group algebra this gives
\(\rho(\operatorname{star}(y))=P^{-1}\rho(y)^tP\).
The identity persists under every commutative scalar extension.
Taking determinants proves the required polarization identity and,
in particular, every \texttt{Pair.P} coefficient identity.
The determinant of \(\rho\) is continuous, and its reduction is
the determinant of the prescribed diagonal tuple.  This proves
\texttt{cont}, \texttt{red} and \texttt{P}.  Distinctness of the
restricted residual characters is the retained hypothesis required
by \texttt{Resid.chib\_inj}; it is not inferred from a generic
base change.

\proofstep{The local characters and ordinary equations.}
Apply AC-G57-integral to the actual map into \(A\).  Its polynomial
projectors give a strict basis over \(A\), so the ordered diagonal
characters \(\Psi_{v,i}\) are continuous \(A^\times\)-valued
characters with the required residual values.  Scalar normalization
multiplies each by \(\overline\kappa_\ell^{-1}\), preserving the
ordered flag.  A continuous character kills the closure of the
commutator subgroup.  On an element whose image in the topological
abelianization has \(\ell\)-power order, its value is therefore
killed by some \(\ell\)-power.  In the reduced characteristic-
\(\ell\) ring \(A\), the equation \(u^{\ell^a}=1\)
implies \((u-1)^{\ell^a}=0\), hence \(u=1\).
This is exactly \texttt{IsPsi.mu} for the stated \texttt{th}.

For any three local elements \(g_0,g_1,g_2\), the ordered product
\[
       \prod_{i=0}^2\bigl(\rho(g_i)-\Psi_{v,i}(g_i)I\bigr)=0
 \tag{G61.6}
\]
in the strict flag basis.  Indeed the rightmost factor maps the
full flag into its second step, the next maps that step into its
first, and the first factor kills the first step.  Conjugation
back to the original basis preserves zero.  Multiplying by any
\(\rho(h)\) and taking traces gives the literal
\texttt{Pair.Ord} equation.  This verifies the complete local
\texttt{psi} and \texttt{Ord} fields, not only the inertia
characteristic polynomial.

\proofstep{Away killing and universality.}
For an away-inertia element in \(\mathfrak r.\texttt{Ip}\), all
three normalized residual characters are one.  The cyclotomic
character is unramified there, so the original residual matrix is
the identity as well.  The actual killing condition gives
\(\rho(g)=I\).  This is the required
\texttt{Pair.I} trace identity.  The stated descent to the frame
ensures the same assertion for its inertia images.
We have now checked every field of \texttt{Pair}.
Its universal property supplies (G61.3) into the actual coefficient
ring \(A\), and its \texttt{Maps} conclusion identifies both the
determinant and all ordered local characters.  Composing with the
injective map \(A\hookrightarrow V\) gives the displayed curve map.
\end{proof}

\begin{proposition}[AC-G61-weight: the actual characteristic-\(\ell\) weight is constant]
\label{AC-G61-weight}
In AC-G61-pair, let \(\Lambda\) be the source
\texttt{LamObj}.  The composite
\[
        \Lambda\xrightarrow{\texttt{toPair}}B
                  \xrightarrow{f}A
 \tag{G61.7}
\]
factors through its residue-field specialization
\(\Lambda\longrightarrow k\hookrightarrow A\).
In particular its image is the finite field \(k\).
\end{proposition}
\begin{proof}
The actual integral local identities prescribe each ordered inertia
character, not merely its semisimplification.  After reduction to
the characteristic-\(\ell\) curve, these are the reductions of
the fixed characters \(\gamma,1,\gamma^{-1}\), followed by
the common normalization \(\overline\kappa_\ell^{-1}\).
Their values lie in \(k^\times\), and they equal the inertia
restrictions of the corresponding residual local characters.
Those residual characters themselves define the constant inertial
datum over \(k\); their \(\ell\)-primary torsion values are one.
Thus this is an \texttt{IsInert} datum and gives a map
\(\Lambda\to k\) by the literal universal property of
\texttt{LamObj}.  Both its composite into \(A\) and (G61.7)
carry every universal inertia character to the actual one.
Uniqueness in that same universal property makes the two maps
equal.  The coefficient map onto \(k\) shows their image is
exactly \(k\).  This argument allows the unramified local
Frobenius characters of the curve to vary; they are not weight
coordinates in \texttt{LamObj}.
\end{proof}

\begin{theorem}[AC-G61: a per-curve Hecke-support criterion for integral M2]
\label{AC-G61}
Retain the exact old \texttt{FinEnd} input (G60.2).
For every actual characteristic-\(\ell\) Type I curve, suppose
that there are data as in AC-G61-pair and an ideal \(J\subset B\)
such that
\[
        B/J\text{ is finite over }\Lambda,
        \qquad J\subseteq\ker(f).
 \tag{G61.8}
\]
The CM field, determinant ring and ideal may depend on the curve.
Then the actual \texttt{M2} holds, including its vertical torsion.
The first assertion of (G61.8) is the literal
\texttt{RingProps.fin} conclusion of the retained generic
\texttt{ModU} supply, when its full domain holds: in particular
\(V.\texttt{S'}=s.\texttt S\), the CM field contains no
primitive \(\ell\)-th root of unity, and the coefficient datum
satisfies \texttt{Large}.  The existence of these particular
host and coefficient choices is retained, not inferred from
the universal determinant definition.  The second is the additional arithmetic
curve-support assertion; it is not a consequence asserted here
of \texttt{ModU}, of A, or of a single classical automorphic seed.
\end{theorem}
\begin{proof}
Since \(J\) is killed, the image of \(B\) in \(A\) is a
quotient of the finite \(\Lambda\)-module \(B/J\).
AC-G61-weight makes it a finite-dimensional \(k\)-algebra.
It is a domain because it is a subring of \(A\), and hence it
is a finite field.  Consequently every coefficient of the
characteristic polynomial of \(\rho(g)\), for every
\(g\in G_L\), is algebraic over \(k\).
Removing the normalization in (G61.2) multiplies coefficients
by powers of an element of \(\mathbb F_\ell^\times\), so
the same conclusion holds for the original representation.

For an arbitrary \(h\in H\), some positive power \(h^a\)
lies in the finite-index subgroup \(G_L\).
Every eigenvalue \(z\) of \(r_\chi(h)\) in an algebraic
closure satisfies that \(z^a\) is an eigenvalue of
\(r_\chi(h^a)\).  It is therefore algebraic over \(k\),
and so is \(z\), even if \(\ell\mid a\).
All characteristic-polynomial coefficients at \(h\) are
algebraic over \(k\).  The elements of \(V=k'[[t]]\)
algebraic over its finite coefficient field belong to \(k'\):
an algebraic element lies in a finite field, and Hensel uniqueness
for \(X^{\ell^b}-X\) identifies it with its constant residue.
The actual residual condition now identifies each coefficient
with that of (G61.1).  The image-algebra argument of AC-G59
excludes an absolutely irreducible rank-three Type I representation
with all those prescribed seed polynomials.

There are therefore no actual Type I curves.  AC-G60-M2, with
precisely its retained old \texttt{FinEnd} hypothesis, yields
the integral \texttt{M2}.  Its conclusion includes vertical
torsion because it uses the exact characteristic-\(\ell\)
DVR criterion, not only a generic-fiber dimension bound.
\end{proof}

\paragraph{The exact source clause still needed.}
\path{WeilGroupCharacters/UniversalOrdinaryDeterminantRing.lean}
defines \texttt{Resid}, \texttt{Pair},
\texttt{UniversalOrdinaryDeterminantRing.univ},
\texttt{LamObj.univ}, \texttt{P1core} and \texttt{P3core}.
The arbitrary residual tuple is allowed in these definitions;
no single-character power condition is imposed.
\path{WeilGroupCharacters/ArithmeticPointClassicality.lean}
defines \texttt{RingProps.fin}, \texttt{ModU} and
\texttt{ProModularRingExists}.
Its \texttt{PT} quantifies over
\(B\to\overline{\mathbb Q}_\ell\) points only.  It does
not apply to the actual map \(B\to V\) of characteristic
\(\ell\).  \texttt{P1core} already assumes
\(J\subseteq\ker f\); \texttt{P3core} requires a starting
stable prime containing \(J\), connected through its stated
sequence of nice reducible points.  Neither begins such a
propagation chain from the existence of the curve alone.

The source conclusion of the required kind is
\texttt{Setup.GoodPt} in
\path{CurveDVRSpecialization/SecondThirdPrimeConditionTransferAlt.lean}:
it produces a map killing \(J\) for its quantified irreducible
characteristic-\(\ell\) lift.  Its \texttt{AbsMain} supply
retains \texttt{Setup}, including \texttt{TwLink},
\texttt{Fr.Ok}, \texttt{Big}, the rank-two input and the
seed/propagation hypotheses.  \texttt{TwLink.chib} imposes
the single-character-power residual tuple, which is not the
cubic tuple (G61.1).  Thus it is not an already proved source
of (G61.8) for this curve.  Establishing generic ordinary
pro-modularity on the actual cubic curve, or supplying a
verified propagation path satisfying the generic \texttt{P1core}
and \texttt{P3core} hypotheses, is the remaining arithmetic
requirement.  No published theorem with all these actual
hypotheses has been verified here.

For clarity, the elementary facts ``finite over weight'' and
``there is a classical point in the support'' alone do not give
the missing containment.  In the commutative-algebra fragment
\(B=\mathcal O[[u,v]]\), \(\Lambda=\mathcal O[[u]]\),
\(J=(v)\), the quotient \(B/J=\Lambda\) is finite free and
the characteristic-zero point \(u=v=0\) kills \(J\).
The fixed-weight curve \(u\mapsto0,\ v\mapsto t\) into
\(k[[t]]\) does not.  This is only a counterexample to that
fragmentary inference, not a countermodel to the full Hida
or propagation statements.

\paragraph{Geometric dictionary and scope.}
The representation and its determinant are the lisse-system and
polynomial-law descriptions of the same compatible finite-level cover data,
in the already audited SGA~1, Expos\'e~V, \S\S4 and 7
framework.  The ideal \(J\) cuts out the ordinary Hecke support
inside the deformation parameter space.  AC-G61-pair and
AC-G61-weight verify the map and weight on the actual curve.
Neither the categorical dictionary nor fixed weight implies
that the curve lies in this support.  The theorem is an exact
sufficient criterion, not a proof of the remaining pro-modularity
input or of a numerical Selmer estimate.
